The current prototype is strongest as a research scaffold: it exposes the full chain from OSM public GPS signals to route reconstruction, feature extraction, profile construction, candidate generation, ranking, and evaluation. The GUI/demo angle remains valuable because it shows how a user would interact with ranked routes and profile explanations.

The main limitation is data validity. OSM public GPS trackpoints cannot be assumed to represent clean, persistent user histories. The current pseudo-history profiles are useful for exploring whether trace-derived movement signals can inform route ranking, but they are not equivalent to logged user route choices. This affects both profile construction and evaluation.

The second limitation is reconstruction quality. GPS-to-route distance is low, but route-distance ratio is high for several segments. This suggests the approximate nearest-node plus shortest-path stitching method can overbuild routes. Before submission, the pipeline should improve map matching, segment splitting, or simplification so strict usability passes without relaxing the route-distance-ratio threshold.

The third limitation is evaluation size. The current route-candidate evaluation uses only three held-out queries. This is enough to expose the weakness of the old history-record ranking setup, but not enough for a publishable claim. The next version should add more cities or bounding boxes, more OSM trace pages where allowed, and preferably a clean trajectory dataset for comparison.

\textbf{TODO:} Add an error analysis table showing which segments fail strict route-distance-ratio checks and why.
